% !TeX root = ../traffic_rules.tex

\section{Traffic rule formalization}
\label{sec:formalization}
We formalize selected rules from the \stvo, e.g., establishing an emergency lane or keeping a safe distance, which are necessary to drive on German interstates and also provide their textual concretization as well as additional comments. The temporal logic formula and the related laws and traffic signs for each rule are listed in Tab.~\ref{tab:overview_traffic_rules}\footnote{The terms \textit{OLG} and \textit{Bay VRS}  are abbreviations for German types of courts.}. Predicates and functions required for the FLTL formulas are introduced in Sec.~\ref{sec:predicates}. We first present general rules and afterward interstate-specific rules. The rules are written from the viewpoint of the autonomous ego vehicle.

\begin{table*}[!tb]\centering
	\caption{Overview about the formalized traffic rules. For brevity we do not show the parameters of the predicates.}
	\center
	\arstretch{1.35}
	\begin{tabular}{@{}lcl@{}} \toprule
		\textbf{Rule name} & \multicolumn{1}{c}{\textbf{Law reference}}  & \multicolumn{1}{c}{\textbf{FLTL formula}}  \\
		\\ 	
		\textbf{R\_G$\mathbf{2}$} & \stvo \S 4~(1); VCoRT \S 17~(1) &  
		\multicolumn{1}{l}{G$(\neg \unnecessaryBraking)$} 
		\\ 
		\textbf{R\_G$\mathbf{3}$} & \multicolumn{1}{c}{\stvo \S 3~(1), \S 3~(3), \S 18~(1), \S 18~(5),} & \multicolumn{1}{l}{G$(\laneSpeedLimitPredicate  \land \fovSpeedLimitPredicate \land  \brakingSpeedLimitPredicate$}\\  
		& \S 18~(6), traffic sign $274$ & $\quad \land \roadConditionSpeedLimitPredicate \land \vehicleTypeSpeedLimitPredicate) $
		\\ 
		\textbf{R\_G$\mathbf{4}$} & \stvo traffic sign $275$ & 
		\multicolumn{1}{l}{G$(\minSpeedLimitPredicate$)} 
		\\
		\textbf{R\_G$\mathbf{5}$} & \stvo \S 3~(2)& 
		\multicolumn{1}{l}{G$(\preservesTrafficFlowPredicate$)} 
		\\
		\textbf{R\_I$\mathbf{1}$} & \multicolumn{1}{c}{ \stvo \S 12 (1), \S 18 (8);} & \multicolumn{1}{l}{G$(\neg(\inCongestion \lor \, \leadVehicleStandingPredicate) \implies \neg \inStandStillPredicate)$} \\
		& OLG Karlsruhe DAR 02, 34; & \\
		& OLG Frankfurt DAR 01, 504; & \\
		& OLG Hamm VersR 91, 83; & \\
		& Bay VRS 59, 54 = StVE 23 & 
		\\
		\textbf{R\_I$\mathbf{2}$} & \multicolumn{1}{c}{ \stvo \S 7 (2), \S 7 (2a), \S 7a} & \multicolumn{1}{l}{G$((\congestionLeftPredicate \land \drivesMaxWithSlightlyHigherSpeedPredicate)$} \\
		&  & $\quad\lor (\rightOfBroadLaneMarkingPredicate \land \leftOfBroadLaneMarkingPredicate)$\\
		&  & $\quad\lor (\onAccessRampPredicate \land \neg \congestionLeftPredicate) \lor \neg \drivesFasterThanVehicleLeftPredicate)$\\
		\textbf{R\_I$\mathbf{3}$} & \stvo \S 18~(7) &  
		G$(\neg \makesUTurnPredicate \land \neg \reversesPredicate)$
		\\
		\textbf{R\_I$\mathbf{4}$} & \stvo \S 11~(2) & 
		\multicolumn{1}{l}{G$(\inCongestion \implies ((\interstateBroadEnoughForEmergencyLanePredicate \implies$} \\ 
		& & $((\inLeftMostLanePredicate \implies \drivesLeftMostPredicate) \land (\neg \inLeftMostLanePredicate \implies \drivesRightMostPredicate)))$ \\
		& & $  \land \, (\neg \interstateBroadEnoughForEmergencyLanePredicate \implies \drivesRightMostPredicate)))$ 
		\\
		\textbf{R\_I$\mathbf{5}$} & \stvo \S 1~(2), \S 18 recital 11 & 
		\multicolumn{1}{l}{G$(\onAccessRampPredicate \land$ F$(\onMainCarriageWayPredicate)$} \\ 
		&  & $ \quad \implies \neg (\neg \inRightMostLanePredicate \land $ F$(\inRightMostLanePredicate)))$
		\\
		\bottomrule
	\end{tabular}
	\label{tab:overview_traffic_rules}
\end{table*} 


\subsection{General traffic rules}
\subsubsection{R\_G1 - Safe distance to preceding vehicle} 
\paragraph{Law reference:} 
\begin{itemize}
	\item \textbf{\stvo \S 4~(1) - English:} A person operating a vehicle moving behind another vehicle must, as a rule, keep a sufficient distance from that other vehicle so as to be able to pull up safely even if it suddenly slows down or stops.
	\item \textbf{\stvo \S 4~(1) - German:} Der Abstand zu einem vorausfahrenden Fahrzeug muss in der Regel so groß sein, dass auch dann hinter diesem gehalten werden kann, wenn es plötzlich gebremst wird.
	\item \textbf{\vcort \S 13~(5):} The  driver  of  a  vehicle  moving  behind  another  vehicle  shall  keep  at a  sufficient distance from that other vehicle to avoid collision if the vehicle in front should suddenly slow down or stop. 
\end{itemize}

\paragraph{Concretization:} The ego vehicle following vehicles within the same lane has to keep a safe distance which ensures collision freedom to each of these vehicles, even if one or several of them are suddenly stopping.

\paragraph{LTL formula:} G$(\inSameLaneAsEgoVehiclePredicate \land \inFrontOfEgoVehiclePredicate \implies \safeDistanceFrontPredicate)$

\paragraph{Comment:}
The rule has to be evaluated for every other traffic participant within the sensor range. Vehicles outside are taken care of by the allowed velocity in rule R\_G3.

\subsubsection{R\_G2 - Unnecessary braking}
\paragraph{Law reference:}
\begin{itemize}
	\item \textbf{\stvo \S 4~(1) - English:} The person operating the vehicle in front must not brake suddenly without a compelling reason.
	\item \textbf{\stvo \S 4~(1) - German:} Wer vorausfährt, darf nicht ohne zwingenden Grund stark bremsen.
	\item \textbf{\vcort \S 17~(1):} No driver of a vehicle shall brake abruptly unless it is necessary to do so for safety reasons.
\end{itemize}


\paragraph{Concretization:} The ego vehicle is not allowed to brake abruptly without reason. % unless it has a reason to do so. 
The ego vehicle brakes abruptly if: a) The ego vehicle brakes although there is no leading vehicle and there is no reason caused by a traffic regulation to reduce the velocity. b) There is an obstacle in front of the vehicle and the acceleration difference and jerk difference do not violate predefined thresholds. c) The maximum allowed speed will change or has changed due to unrecognized traffic signs or environmental conditions. All possible braking situations which are not covered by the three cases are unnecessary.

\paragraph{LTL formula:} G$(\neg \unnecessaryBraking)$

\subsubsection{R\_G3 - Maximum speed limit}
\paragraph{Law reference:}
\paragraph{Concretization:} The ego vehicle is not allowed to exceed the speed limit of the lanes it drives on, the speed limit which is necessary to ensure collision freedom with respect to vehicles outside the field of view, the speed which is necessary to comfortably react to traffic regulations, the speed limit based on road conditions, and the maximum velocity which is allowed for its vehicle type.

\subsubsection{R\_G4 - Minimum speed limit}
\paragraph{Law reference:}
\paragraph{Concretization:} The ego vehicle has to consider the minimum speed limit traffic sign.

\subsubsection{R\_G5 - Traffic flow}
\paragraph{Law reference:}
\paragraph{Concretization:} The ego vehicle is not allowed to travel so slowly that it impedes traffic flow.

\subsection{Interstate-specific traffic rules}

\subsubsection{R\_I1 - Stopping}
\paragraph{Law reference:}
\paragraph{Concretization:} If the ego vehicle is not part of a congestion or the directly leading vehicle is not standing, then stopping on the main carriage way, shoulder lane, and ramps is forbidden.

We allow the ego vehicle to stop if a leading vehicle also stops.

\subsubsection{R\_I2 - Driving right faster than left traffic}
\paragraph{Law reference:}
\paragraph{Concretization:} The ego vehicle is not allowed to drive faster than any vehicle on lanes left of it. An exception is, if a) the vehicle in a left lane is in a congestion; b) the vehicle in the left lane is in a congestion on the main carriage way and the ego vehicle is on an off ramp; c) the left vehicle and the ego vehicle are separated by a broad lane marking;  d) the ego vehicle is within an access ramp. For a) and b), the ego vehicle is allowed to drive faster than the vehicle in the left lane with only a slightly higher speed. 

\paragraph{Comment:} The second case of R\_I2 is covered in the first clause of the corresponding FLTL formula in Tab~\ref{tab:overview_traffic_rules}. The fourth case of R\_I2 allows to drive on access ramps with higher speed than on the main carriage way. However, if a congestion is on the left side of the ego vehicle and the ego vehicle is on an access ramp, it is only allowed to drive with slightly higher speed. This rule has to be evaluated with respect to every other vehicle within the sensor range.

\subsubsection{R\_I3 - No reversing and turning}
\paragraph{Law reference:}
\paragraph{Concretization:} Making U-turns and reversing is prohibited.

\subsubsection{R\_I4 - Emergency lane}
\paragraph{Law reference:}
\paragraph{Concretization:} A vehicle in a congestion has to drive leftmost if it is in the leftmost lane and has to drive rightmost otherwise. If the interstate is not broad enough to create an emergency lane, all vehicles should drive in the rightmost lane.

\subsubsection{R\_I5 - Consider entering vehicles}
\paragraph{Law reference:}
\paragraph{Concretization:} The ego vehicle is not allowed to move from the main carriage way to the rightmost lane if another vehicle is about to enter the main carriage way from the access ramp.





%
%
%
%\section{Traffic Rule Formalization}
%\label{sec:rule_formalization}
%Rules are formalized to account for general situations, e.g. the handling after collisions is not formalized.
%\subsection{Traffic Law Terminology}
%\label{sec:def_law_terms}
%\begin{itemize}
%	\item access ramp / acceleration lane
%	\item off ramp / deceleration lane
%	\item main carriage way
%	\item motor road / motor street / highway
%	\item ego vehicle
%	\item other vehicle
%\end{itemize}
%
%\subsection{General Traffic Rules}
%\label{sec:general_rules}
%
%\subsubsection{Position}

%\begin{table}[H]
%	\centering
%	\begin{tabular}[t]{ll}
%		\toprule
%		\textbf{Lane following} & \\
%		\textbf{vertical position} & \\
%		\midrule
%		StVO \S 2 (2)& Es ist möglichst weit rechts zu fahren, nicht nur bei Gegenverkehr, \\
%		& beim Überholtwerden, an Kuppen, in Kurven oder bei Unübersichtlichkeit. \\ 
%		& Drivers and riders must always keep as far to the right as possible, not only when \\
%		& traffic is approaching from the opposite direction, when being overtaken, \\
%		& when approaching the brow of a hill, on bends or when their view ahead is obstructed.\\
%		&\\
%		VCoRT \S 12 (1) & When passing oncoming traffic, a driver shall leave sufficient lateral space and, \\
%		& if necessary, move close to the edge of the carriageway appropriate to the direction of traffic. \\
%		& If in so doing he finds his progress impeded by an obstruction or by the presence of other \\ 
%		& road-users, he shall slow down and if necessary stop to allow the oncoming road-user \\
%		& or road-users to pass. \\
%		&\\
%		Concretization &  If there exist only one lane on a road and vehicles can drive in both directions on this lane, \\
%		& the ego vehicle should drive as to the right as possible. If there exist at least one lane \\
%		& for each driving direction and the ego vehicle's lane is at least as broad as the ego vehicle \\
%		& or there is no adjacent lane in opposite driving direction, the ego vehicle should \\
%		& drive in the lane center.\\
%		&\\
%		LTL & G$((\onlyOnLanePredicate \land \noLaneInOppositeDirectionPredicate \implies \drivesFarRightInLanePredicate)   $\\
%		& $ \land (\twoLanesWithDifferentDirectionsPredicate \land \egoLaneBroadEnoughPredicate$ \\
%		& $\implies \drivesOnLaneCenterPredicate))$\\
%		MTL & G$(\safeDistanceFrontPredicate)$\\
%		STL & G$(\safeDistance(\vEgo, \vOther) \leq (\rearPositionFunction(\otherVehicle) - \frontPositionFunction(\egoVehicle)))$\\
%		\bottomrule
%	\end{tabular}
%\end{table}
%
%
%\subsubsection{Braking}
%\begin{table}[H]
%\centering
%\begin{tabular}[t]{ll}
%\toprule
%\textbf{Unnecessary braking} & \\
%\midrule
%	StVO \S 4 (1) & Wer vorausfährt, darf nicht ohne zwingenden Grund stark bremsen.\\
%	& The person operating the vehicle in front must not brake suddenly \\
%	& without a compelling reason.\\
%	&\\
%	VCoRT \S 17 (1) &  No driver of a vehicle shall brake abruptly unless it is necessary to do so for\\
%	&  safety reasons.\\
%	&\\
%	Concretization &  A vehicle shall never brake abruptly unless it has a reason to do so. \\
%	& A vehicle brakes abruptly if: \\
%	& 1. The vehicle brakes although there is no leading vehicle and there is no reason caused by \\
%	& a traffic regulation to reduce the vehicle's velocity; \\ 
%	& 2. There is an obstacle in front of the vehicle and the acceleration difference and \\
%	& jerk difference between the vehicle and the obstacle violate thresholds; \\
%	& 3. The maximum allowed speed will change or changed due to traffic \\
%	& signs or environmental conditions.\\
%	&\\
%	LTL & G$(\unnecessaryBraking)$\\
%	MTL & G$(\unnecessaryBraking)$\\
%	STL & \\
%\bottomrule
%\end{tabular}
%\end{table}%
%
%\begin{table}[H]
%	\centering
%	\begin{tabular}[t]{ll}
%		\toprule
%		\textbf{No stopping} & \\
%		\midrule
%		StVO \S 18 (8), \S 12 (1), & Halten, auch auf Seitenstreifen, ist verboten.\\
%		OLG Karlsruhe DAR 02, 34,& \\
%		OLG Frankfurt DAR 01, 504, & Das Halten ist unzulässig auf Einfädelungs- und auf Ausfädelungsstreifen.\\
%		OLG Hamm VersR 91, 83, & \\
%		Bay VRS 59, 54 = StVE 23 & Das Anhalten auf allen Teilen der Autobahn, die dem Zweck dienen einen \\
%		 & flüssigen Verkehr zu gewährleisten, wie die Autobahnen selbst, Zufahrtswege \\
%		 & oder dem Seitenstreifen, ist verboten. \\
%		StVO English & \\
%		&\\
%		VCoRT \S 17 (1) &  \\
%		&\\
%		Concretization & Stopping on the main carriage way, shoulder lane, and ramps is forbidden.  \\
%		&\\
%		LTL & G$(\neg \inCongestion \land (\onAccessRampPredicate \lor \onExitRampPredicate $ \\
%		& $\lor \onMainCarriageWayPredicate \lor \onShoulderPredicate) \implies \neg \inStandStillPredicate)$\\
%		MTL & G$(\neg \inCongestion \land (\onAccessRampPredicate \lor \onExitRampPredicate $ \\
%		& $\lor \onMainCarriageWayPredicate \lor \onShoulderPredicate) \implies \neg \inStandStillPredicate)$\\
%		STL & \\
%		\bottomrule
%	\end{tabular}
%\end{table}%
%
%\subsubsection{Velocity}
%\begin{table}[H]
%\centering
%\begin{tabular}[t]{ll}
%\toprule
%\textbf{Maximum Speed Limit} &\\
%\midrule
%	StVO \S 3 (1), \S 3 (3),  & The legal text is omitted for reasons of space.\\
%	\S 18 (1), \S 18 (5), \S 18 (6) & \\
%	&\\
%	VCoRT \S 13 (2) &  Domestic legislation shall establish maximum speed limits for all roads. \\
%	& Domestic legislation shall also determine special speed limits applicable \\ 
%	& to certain categories of vehicles presenting a special danger, in particular by reason \\ & of their mass or their  load. They may establish similar provisions \\
%	& for certain categories of drivers, in particular for new drivers.\\
%	&\\
%	Concretization &  A vehicle shall not exceed the allowed speed limit of the lanes it drives on, shall \\
%	& not exceed the speed limit which is necessary to ensure collision freedom with \\
%	& respect to vehicles outside the field of view, and shall not exceed as speed which \\
%	& is necessary to comfortably ensure traffic regulations.\\
%	&\\
%	Comments &  The traffic laws do not provide concrete information about speed limit, \\ 
%	& e.g., the VCoRT passes the task of defining speed limits to national lawmakers. \\
%	& The concretization assumes that a mapping of speed limits provided \\ 
%	& by traffic signs to lanes is given. If no speed limit is given, \\ 
%	& the speed limit caused by the field of view (traffic sign based speed limit set \\ 
%	& to infinity) and a speed limit which ensures comfortable braking for traffic regulations \\
%	& (e.g., upcoming speed limit) must be considered. The velocity with respect to vehicles \\
%	& inside the field of view is considered in the safe distance traffic rule concretization.\\
%	&\\
%	LTL & G$(\laneSpeedLimitPredicate \land \fovSpeedLimitPredicate \land$ \\
%	& $\brakingSpeedLimitPredicate \land \roadConditionSpeedLimitPredicate)$\\
%	MTL & G$(\laneSpeedLimitPredicate \land \fovSpeedLimitPredicate \land$\\
%	& $\brakingSpeedLimitPredicate \land \roadConditionSpeedLimitPredicate)$\\
%	STL & G($\vEgo \leq \speedLimitLane \land \vEgo \leq \speedLimitFOV \land \vEgo \leq \speedLimitBraking \land \vEgo \leq \speedLimitWeather)$\\ 
%\bottomrule
%\end{tabular}
%\end{table}
%
%\begin{table}[H]
%	\centering
%	\begin{tabular}[t]{ll}
%		\toprule
%		\textbf{Minimum Speed Limit} &\\
%		\midrule
%		StVO \S 3 (2)  & Ohne triftigen Grund dürfen Kraftfahrzeuge nicht so langsam fahren, dass sie\\
%		& den Verkehrsfluss behindern.\\
%		& No motor vehicle must, without good reason, travel so slowly as to impede\\ 
%		& the flow of traffic.\\
%		sign $275$ &  Wer ein Fahrzeug führt, darf nicht langsamer als mit der angegebenen \\ 
%		& Mindestgeschwindigkeit fahren, sofern nicht Straßen-, Verkehrs-, Sicht- oder \\ 
%		& Wetterverhältnisse dazu verpflichten. Es verbietet, mit Fahrzeugen, die nicht so \\
%		& schnell fahren können oder dürfen, einen so gekennzeichneten Fahrstreifen zu benutzen. \\
%		& A person operating a vehicle must not do so below the indicated minimum speed \\ 
%		& unless road, traffic, visibility or weather conditions make this necessary.  \\
%		& This sign prohibits vehicles that are unable or not permitted to travel as fast from \\
%		& using a lane marked in this manner. \\
%		&\\
%		VCoRT \S 13 (2) &  Similar to maximum speed limit.\\
%		&\\
%		Concretization & A vehicle shall consider the minimum speed limit based on a corresponding \\
%		& traffic sign and should not travel so slowly to impede traffic flow. \\
%		&\\
%		LTL & G$(\minSpeedLimitPredicate \land \preservesTrafficFlowPredicate)$\\
%		MTL & G$(\minSpeedLimitPredicate \land \preservesTrafficFlowPredicate)$\\
%		STL & G$()$\\
%		\bottomrule
%	\end{tabular}
%\end{table}
%
%
%\subsection{Highway Traffic Rules}
%\label{sec:highway_rules}
%\begin{table}[H]
%	\centering
%	\begin{tabular}[t]{ll}
%		\toprule
%		\textbf{Driving right} & \\
%		\textbf{faster than left} &\\
%		\midrule
%		StVO \S 7 (2), \S 7 (2a),  & Ist der Verkehr so dicht, dass sich auf den Fahrstreifen für eine Richtung \\
%		\S 7a & Fahrzeugschlangen gebildet haben, darf rechts schneller als links gefahren werden.\\
%		& \\
%		& Wenn auf der Fahrbahn für eine Richtung eine Fahrzeugschlange auf dem jeweils linken \\ 
%		& Fahrstreifen steht oder langsam fährt, dürfen Fahrzeuge diese mit geringfügig höherer \\ 
%		& Geschwindigkeit und mit äußerster Vorsicht rechts überholen. \\ 
%		& \\
%		& Auf Ausfädelungsstreifen darf nicht schneller gefahren werden als auf den \\  
%		& durchgehenden Fahrstreifen. Stockt oder steht der Verkehr auf den durchgehenden \\ 
%		& Fahrstreifen, darf auf dem Ausfädelungsstreifen mit mäßiger Geschwindigkeit  \\
%		& und besonderer Vorsicht überholt werden. \\
%		& \\
%		& Gehen Fahrstreifen, insbesondere auf Autobahnen und Kraftfahrstraßen, \\
%		& von der durchgehenden Fahrbahn ab, darf beim Abbiegen vom Beginn einer breiten \\
%		& Leitlinie (Zeichen 340) rechts von dieser schneller als auf der durchgehenden Fahrbahn \\
%		& gefahren werden. \\ 
%		& \\
%		& Auf Ausfädelungsstreifen darf nicht schneller gefahren werden als auf den durchgehenden \\ 
%		& Fahrstreifen. Stockt oder steht der Verkehr auf den durchgehenden Fahrstreifen, darf \\ 
%		& auf dem Ausfädelungsstreifen mit mäßiger Geschwindigkeit und besonderer Vorsicht \\
%		& überholt werden. \\
%		
%		StVO English & If the density of traffic has resulted in queues of vehicles in the lanes of one \\
%		& carriageway, traffic on the right (nearside lane, middle lane) may move faster than \\
%		& traffic on the left (offside lane, middle lane).\\
%		& \\
%		& If, on a carriageway for one direction of traffic, a vehicle queue has formed and come to \\ 
%		& a standstill or is moving at low speed in the left-hand lane, vehicles may overtake this \\ 
%		& queue on the right at a slightly higher speed and with the utmost care. \\
%		& \\
%		& In deceleration lanes, vehicles must not travel faster than traffic on the main carriageway. \\
%		& If traffic on the main carriageway is moving slowly or is stationary, vehicles in a \\
%		& deceleration lane may overtake at a moderate speed and with great care. \\
%		& \\
%		& Where lanes diverge from the main carriageway, especially on motorways and motor roads, \\
%		& vehicles turning off may, from the beginning of and to the right of a lane marking with \\
%		& broad lines (sign 340), travel faster than traffic on the main carriageway. \\
%		& \\
%		& On motorways and other roads outside built-up areas, vehicles may travel faster in \\
%		& acceleration lanes than traffic on the main carriageway. \\
%	\end{tabular}
%\end{table}
%
%\begin{table}[H]
%	\centering
%	\begin{tabular}[t]{ll}
%		\toprule
%		\textbf{Driving right} & \\
%		\textbf{faster than left} &\\
%		\midrule
%		VCoRT  &  \\
%		
%		Concretization &  A vehicle shall not drive faster than any vehicle on its left lanes. \\
%		& A exception is, if $1.$ the vehicle on the left lane is in a congestion; \\
%		& $2.$ the vehicle on the left lane is on the main carriage way and the vehicle on the right \\
%		& lane is on a off ramp; $3.$ the left vehicle and right vehicle are separated by a broad \\
%		& lane marking;  $4.$ the right vehicle is within an access ramp. For $1.$ and $2.$, the vehicle \\ 
%		& on the right lane is allowed to drive faster than the vehicle on the left lane \\
%		& with only a slightly higher speed.\\
%		
%		Predicates & $\drivesFasterThanVehicleLeftPredicate$, $\congestionLeftPredicate$, $\otherVehicleLeftOfBroadLaneMarkingPredicate$,\\
%		& $\drivesMaxWithSlightlyHigherSpeedPredicate$, $\rightOfBroadLaneMarkingPredicate$\\
%		
%		Comments &  Assumption: Deceleration lane always right of main carriage way.\\
%		& Case two can be neglected -> same as case one \\
%		
%		LTL & G$(\neg \drivesFasterThanVehicleLeftPredicate$ \\ 
%		& $\lor (\congestionLeftPredicate \land \drivesMaxWithSlightlyHigherSpeedPredicate)$\\
%		& $\lor (\rightOfBroadLaneMarkingPredicate \land \otherVehicleLeftOfBroadLaneMarkingPredicate)$\\
%		& $\lor (\onAccessRampPredicate \land \neg \congestionLeftPredicate))$\\
%		MTL & G$(\neg \drivesFasterThanVehicleLeftPredicate$ \\ 
%		& $\lor (\congestionLeftPredicate \land \drivesMaxWithSlightlyHigherSpeedPredicate))$\\
%		STL & \\
%		\bottomrule
%	\end{tabular}
%\end{table}
%
%
%\begin{table}[H]
%	\centering
%	\begin{tabular}[t]{ll}
%		\toprule
%		\textbf{Emergency Lane} &\\
%		\midrule
%		StVO \S 3 (2) & Sobald Fahrzeuge auf Autobahnen sowie auf Außerortsstraßen mit mindestens zwei \\
%		& Fahrstreifen für eine Richtung mit Schrittgeschwindigkeit fahren oder sich die Fahrzeuge \\
%		& im Stillstand befinden, müssen diese Fahrzeuge für die Durchfahrt von Polizei- \\
%		& und Hilfsfahrzeugen zwischen dem äußerst linken und dem unmittelbar rechts \\
%		& daneben liegenden Fahrstreifen für eine Richtung eine freie Gasse bilden.\\
%		
%		StVO English & As soon as vehicles on motorways and roads outside a built-up area with at least two lanes \\
%		& for one direction start to move at walking pace or come to a standstill, these vehicles must \\
%		& leave a gap for one direction between the lane on the far left and the lane immediately \\
%		& adjacent to it on the right to allow police and emergency vehicles to pass.\\
%		
%		VCoRT  &  \\
%		
%		Concretization &  A vehicle in a congestion outside urban roads has to drive leftmost if it is in the leftmost \\
%		& lane and has to drive rightmost otherwise. \\
%		
%		Predicates & $\inLeftMostLanePredicate$, $\drivesLeftMostPredicate$\\
%		
%		Comments &  \\
%		
%		LTL & G$((\inLeftMostLanePredicate \implies \drivesLeftMostPredicate)$ \\
%		& $\land (\neg \inLeftMostLanePredicate \implies \drivesRightMostPredicate))$\\
%		MTL & G$((\inLeftMostLanePredicate \implies \drivesLeftMostPredicate)$ \\
%		& $\land (\neg \inLeftMostLanePredicate \implies \drivesRightMostPredicate))$\\
%		STL & \\
%		\bottomrule
%	\end{tabular}
%\end{table}
%
%\begin{table}[H]
%	\centering
%	\begin{tabular}[t]{ll}
%		\toprule
%		\textbf{No reversing}&\\ 
%		\textbf{and turning} &\\
%		\midrule
%		StVO \S 18 (7) & Wenden und Rückwärtsfahren sind verboten.\\
%		
%		StVO English & Making U-turns and reversing are prohibited.\\
%		
%		VCoRT  &  \\
%		
%		Concretization &  Making U-turns and reversing are prohibited.\\
%		
%		Predicates & $\makesUTurnPredicate$, $\reversesPredicate$\\
%		
%		Comments &  \\
%		
%		LTL & G$(\neg \makesUTurnPredicate \land \neg \reversesPredicate)$ \\
%		MTL & G$(\neg \makesUTurnPredicate \land \neg \reversesPredicate)$\\
%		STL & \\
%		\bottomrule
%	\end{tabular}
%\end{table}
%
%\begin{table}[H]
%	\centering
%	\begin{tabular}[t]{ll}
%		\toprule
%		\textbf{Direction indicator}&\\ 
%		\midrule
%		StVO \S 5 (4a) & Das Ausscheren zum Überholen und das Wiedereinordnen sind rechtzeitig und \\
%		& deutlich  anzukündigen; dabei sind die Fahrtrichtungsanzeiger zu benutzen. \\
%		& The intention of pulling out to overtake and of moving back to the right-hand side \\
%		& of the road are to be signalled clearly and in good time, using the vehicle’s \\
%		& direction indicators to do so.\\
%		&\\
%		StVO \S 7 (5) & [...] Jeder Fahrstreifenwechsel ist rechtzeitig und deutlich anzukündigen; \\
%		& dabei sind die Fahrtrichtungsanzeiger zu benutzen. \\
%		& [...] Any change of lane must be signalized clearly and in good time, \\
%		& using the vehicle’s direction indicators to do so. \\
%		&\\
%		Burmann, et al., &  Der Fahrtrichtungsanzeiger is rechtzeitig bevor Verlassen der Hauptfahrbahn und\\
%		Straßenverkehrsrecht,&  dem Befahren des Ausfädelungsstreifen zu benutzen. Sobald der Ausfädelungsstreifen\\
%		Recital 11  &  befahren wird, soll der Fahrtrichtungsanzeiger deaktiviert werden. \\
%		
%		& The direction indicator shall be timely switched on before leaving \\
%		& the main carriageway and moving to the exit ramp. When the \\
%		& exit ramp had been entered, the direction indicator should be switched off. \\
%		&\\
%		VCoRT  &  \\
%		&\\
%		Concretization &  Any change of lane must be signalized clearly and enough time before \\
%		& entering the new lane, using the vehicle’s direction indicators.\\
%		& When the new lane had been entered, the direction indicator should be switched off.\\
%		&\\
%		Predicates & \\
%		&\\
%		Comments & We assume the exit ramp is long enough so that a vehicle can switch off the direction indicator before the vehicle has to turn the indicator on for a possible intersection.  \\
%		&\\
%		LTL & G$()$ \\
%		MTL & G$()$\\
%		STL & \\
%		\bottomrule
%	\end{tabular}
%\end{table}
%
%\subsection{Urban Traffic Rules}
%\label{sec:urban_rules}
%
%\subsection{Cross-Country Road Traffic Rules}
%\label{sec:cross_country_rules}